\section{第\ref{sec:code}章　摩尔斯电码}

本章的极限估计来自信息论和计算；测量过的是信道中噪声产生在哪里、大脑工作得多快、一个脉冲要花多少代价，而皮层实际使用哪种编码，仍是悬而未决的问题。

{\footnotesize
\setlength{\tabcolsep}{3pt}\renewcommand{\arraystretch}{1.15}
\begin{longtable}{>{\raggedright}p{0.5cm}>{\raggedright}p{4.0cm}>{\raggedright}p{5.6cm}>{\raggedright}p{2.3cm}>{\raggedright\arraybackslash}p{1.7cm}}
\toprule
编号 & 命题 & 实验与测量内容 & 来源 & 状态 \\
\midrule\endfirsthead
\toprule
编号 & 命题 & 实验与测量内容 & 来源 & 状态 \\
\midrule\endhead
1 & 皮层\nt{GLUT}神经元的频率从每秒零点几个到几十个脉冲，神经元大部分时间工作在下限附近 & 在清醒大鼠听皮层中用贴附细胞的移液管记录（神经元的选择与其活动无关）：任一时刻频率高于每秒 $20$ 个脉冲的神经元不到 $5\,\%$；部分神经元的背景频率低于每秒 $0.01$；频率分布为对数正态 & \cite{Hromadka2008} & 已测量 \\
2 & 脉冲信道容量 $R_{\max}\approx r\log_2\frac{e}{r\Delta t}$~\eqref{eq:spikeentropy}，几乎全部在于脉冲的时间（命题~\ref{prop:capacity}） & 由长度为 $\Delta t$ 的格子组成的二进制串的熵。本章的数字：$r=10$ 每秒、$\Delta t=1\ms$ 时每个脉冲 $8$ 比特，$\Delta t=10\ms$ 时约 $5$ 比特，脉冲计数器每秒不到 $2$ 比特 & \cite{MacKay1952} & 理论 \\
3 & 感觉神经元每个脉冲传送一到几比特，最好的情况下可达上限的一半 & 对已知刺激的感觉神经元，根据其脉冲重建刺激；测量汇总见该书 & \cite{Rieke1997} & 已测量 \\
4 & 脉冲发生器是精确的；精确时间由输入的快速跳变决定 & 大鼠皮层脑片中的神经元：对重复的、快速波动的电流，脉冲以优于 $1\ms$ 的精度重复；对恒定电流，脉冲时刻漂移开来 & \cite{Mainen1995} & 已测量 \\
5 & 突触不可靠：一半以上的脉冲因释放的随机性而到不了接收者 & 对海马锥体神经元的输入做最小刺激（脑片）：一半以上的脉冲不引起响应。排除了刺激阈值波动、轴突分支处的传导故障和实验温度低等因素；剩下的是概率性释放 & \cite{AllenStevens1994} & 已测量 \\
6 & 活体皮层中的脉冲流看起来是随机的：间隔散布如泊松流，脉冲数方差接近均值；部分“噪声”是关于动物运动的消息 & 清醒猴视区V1和MT的记录：间隔的变异系数约为 $1$，脉冲数方差与均值之比如同随机过程；把独立小输入相加的模型给出的散布要小得多。在小鼠视皮层中，相当大一部分散布可以由动物自身运动的录像预测 & \cite{Softky1993,Stringer2019} & 已测量 \\
7 & 频率编码很慢：误差 $1/\sqrt{rT}$~\eqref{eq:rateerror}，而大脑在 $\sim150\ms$ 内认出图片 & 公式是泊松统计，本章推导。实验：受试者判断呈现 $20\ms$ 的陌生照片中是否有动物；“是”和“否”的脑电位约在 $150\ms$ 后分开。把这段时间分成十来级、每级 $10$--$20\ms$，是本章的估计，不是测量 & \cite{Thorpe1996} & 已测量 \\
8 & 对神经元群求平均受相关性限制：误差最多减小到 $1/\sqrt c$ 分之一~\eqref{eq:correlated}（命题~\ref{prop:pooling}） & 同时记录的猴MT区神经元对：脉冲数相关系数平均为 $0.12$，理论分析表明即便这样的相关也限制了神经元群的收益。理论：只有与信号相似的那部分相关设定极限。对立观点的模型：$50$--$100$ 个神经元的群体频率在一个脉冲间隔（$10$--$50\ms$）内就能读出，皮层不利用单个脉冲的时间 & \cite{Zohary1994,MorenoBote2014,ShadlenNewsome1998} & 已测量 \\
9 & 数可以写在首个脉冲的时间~\eqref{eq:latency}、神经元群的脉冲顺序或局部节律的相位中 & 蝾螈视网膜：短暂呈现图像的空间结构记录在输出神经元首个脉冲的相对延迟中，几乎与对比度无关。大鼠海马位置细胞：穿过“自己的”位置时，脉冲在 $\theta$ 节律（$7$--$12$\,Hz）周期中越来越早，移动 $100^\circ$ 到 $355^\circ$。皮层在多大程度上利用脉冲时间，是悬而未决的问题（第~8行） & \cite{Gollisch2008,OKeefe1993} & 已测量 \\
10 & 读出哪种编码由接收者的时间常数决定~\eqref{eq:reader}（命题~\ref{prop:reader}）；活跃皮层中的神经元更接近符合检测器 & 公式~\eqref{eq:reader} 是本章推导。在体记录综述：活跃皮层中膜电导比静息时高数倍，时间常数相应缩短。皮层正是以这种方式切换编码，尚无直接证明 & \cite{Destexhe2003} & 假说 \\
11 & 压抑型突触传送频率的变化，本质上是 $\Delta r/r$~\eqref{eq:depression}（图~\ref{fig:depression}） & 皮层锥体神经元成对记录：压抑速度由释放概率决定；压抑慢时响应反映频率，压抑快时反映符合。基于大鼠皮层实测压抑的模型：快输入和慢输入上相同的频率相对变化给出相同的响应，即突触传送 $\Delta r/r$。数字 $1.7\to2.2$、$6.7$ 和 $90\ms$ 是本章的计算 & \cite{Tsodyks1997,Abbott1997depression} & 已测量 \\
12 & 簇放电是“划”：它丢失的概率 $(1-p)^k$~\eqref{eq:burst} 很小，易化使之更小 & 综述：在不可靠的突触上，单个脉冲常常丢失，而簇放电借助易化得以可靠传送；簇放电引起突触的长时程变化。大脑把簇放电当作独立符号来读，是假说 & \cite{Lisman1997,AllenStevens1994} & 假说 \\
13 & 快速\nt{GABA}神经元设定节律和“标点”；另一类抑制抑制性神经元并打开闸门（命题~\ref{prop:punctuation}） & 皮层\nt{GABA}神经元各类别特性的综述。光遗传学：节律性兴奋快速\nt{GABA}神经元会诱发 $\gamma$ 节律，对刺激的响应取决于周期的相位。在小鼠视皮层中，PV神经元相互抑制，SST神经元抑制其他所有类别，VIP神经元主要抑制SST。在清醒小鼠中兴奋VIP神经元会解除对\nt{GLUT}神经元的抑制；开启VIP神经元的是强化信号。作为一般规则的角色分工是假说 & \cite{Tremblay2016,Cardin2009,Pfeffer2013,Pi2013} & 假说 \\
14 & 能量把平均频率限制在每秒约一个脉冲的量级~\eqref{eq:energy} & 根据解剖和生理数据得出的啮齿类灰质能量预算：脉冲和突触电流分别占信号传递能量的 $47\,\%$ 和 $34\,\%$；同时活跃的神经元不超过 $\sim15\,\%$。人类皮层：同时强烈活跃的神经元不超过 $\sim1\,\%$。每个脉冲 $\sim10^9$ 个ATP、信号传递功率 $\sim10$\,W，是本章基于这些计算的估计 & \cite{Attwell2001,Lennie2003} & 理论 \\
15 & 编码是稀疏的，并被调整到每焦耳比特数最多；皮层以精度换能量（命题~\ref{prop:economy}） & 节约编码假说。依据：限制食物的小鼠视皮层中，AMPA受体电导和突触ATP消耗下降（$-29\,\%$），脉冲频率保持不变，而朝向调谐变宽 $32\,\%$ & \cite{Barlow1961,Padamsey2022} & 假说 \\
\bottomrule
\end{longtable}}
